\setcounter{section}{0}
\renewcommand{\thesection}{\Alph{section}}
\renewcommand{\theHsection}{anexo.\arabic{section}}
\counterwithin{table}{section}
\renewcommand{\theHtable}{anexo.\arabic{section}.\arabic{table}}

\begin{landscape}
\chapter*{ANEXOS}
\addcontentsline{toc}{chapter}{ANEXOS}

\section{MATRIZ DE CONSISTENCIA}
\label{anx:matriz-consistencia}

La matriz conserva la relación entre problemas, objetivos, productos y evidencia. Se conservan HE1--HE4 como proposiciones contrastables del plan; los criterios determinan el estado de sus productos y no sustituyen las hipótesis por una declaración automática de cumplimiento.

\begin{table}[H]
\centering
\caption{Matriz de consistencia de la investigación. Fuente: Elaboración propia a partir del plan de tesis.}
\label{tab:matriz-consistencia}
\scriptsize
\setlength{\tabcolsep}{3pt}
\begin{tabularx}{\linewidth}{L{0.8cm}Y Y Y Y}
\toprule
\textbf{ID} & \textbf{Problema específico} & \textbf{Objetivo específico} & \textbf{Producto} & \textbf{Evidencia y criterio} \\
\midrule
1 & No existe una especificación formal que ordene cable, materiales, fuentes, contornos, interfaces y $\kx$. & Elaborar una especificación trazable de entradas, requisitos y criterios. & Especificación y matriz requisito--prueba. & Cobertura de requisitos, consistencia de unidades y trazabilidad. \\
2 & No se dispone de un paquete integrado de referencias para verificar el artefacto. & Construir y verificar el artefacto con referencias analíticas, manufacturadas, FEM y publicadas. & Artefacto versionado y expediente de aceptación. & Error de $\Tmax$, NRMSE, residuos, balance, convergencia y robustez. \\
3 & No se ha cuantificado de forma trazable el efecto de contraste, patrón, extensión y proximidad. & Aplicar el artefacto a escenarios controlados y consolidar la evidencia térmica. & Expediente de escenarios heterogéneos. & $T(x,y)$, $\Tmax$, $\Delta\Tmax$, punto caliente y cobertura. \\
4 & No se conocen las condiciones de mayor discrepancia entre representaciones heterogéneas y homogéneas. & Comparar la ampacidad de representaciones pareadas e identificar las mayores discrepancias. & Expediente comparativo de ampacidad. & $\Imax$, $\Delta I$, margen térmico, clasificación y límites. \\
\bottomrule
\end{tabularx}
\end{table}
\end{landscape}

Los productos corresponden a la especificación y los expedientes de la formulación explícita. La evidencia anterior de reconstrucción radial se conserva como antecedente. El cumplimiento de OE4 requiere una solución de corriente límite y su comparación FEM.

\section{MATRIZ DE OPERACIONALIZACIÓN}
\label{anx:operacionalizacion}

La operacionalización sigue una cadena de objetos. El producto de una etapa se convierte en parte de la entrada de la siguiente. Los parámetros físicos describen los casos y no se confunden con variables estadísticas de una población.

Se mantiene la cadena del plan: $\mathrm{VI1}=P_1$, $\mathrm{VI2}=\{\mathrm{VD1},R_2\}$, $\mathrm{VI3}=\{\mathrm{VD2},E_3\}$ y $\mathrm{VI4}=\{\mathrm{VD3},C_4\}$. Los objetos nuevos son el paquete físico $P_1$, las referencias $R_2$, los escenarios térmicos $E_3$ y los pares de representación $C_4$. Los indicadores numéricos están contenidos en los expedientes; no reemplazan esos objetos ni convierten las semillas en una población de instalaciones. Un producto restringido no se hereda como aceptación universal en la etapa siguiente.

\begin{table}[H]
\centering
\caption{Operacionalización por etapas del artefacto. Fuente: Elaboración propia.}
\label{tab:operacionalizacion}
\scriptsize
\setlength{\tabcolsep}{3pt}
\begin{tabularx}{\textwidth}{L{1.0cm}Y Y Y}
\toprule
\textbf{Etapa} & \textbf{Entrada} & \textbf{Salida} & \textbf{Indicadores} \\
\midrule
OE1 & Paquete de geometría, materiales, fuentes, contornos y requisitos. & Especificación trazable. & Cobertura, trazabilidad, unidades y supuestos. \\
OE2 & Especificación y referencias independientes. & Artefacto verificado y expediente de aceptación. & Error, NRMSE, residuos, balance, robustez y dominio. \\
OE3 & Artefacto aceptado y catálogo de escenarios. & Expediente de evidencia térmica. & $T(x,y)$, $\Tmax$, $\Delta\Tmax$, gradiente y punto caliente. \\
OE4 & Evidencia térmica y catálogo de pares. & Expediente comparativo de ampacidad. & $\Imax$, $\Delta I$, margen y clasificación del caso. \\
\bottomrule
\end{tabularx}
\end{table}

\section{MATRIZ REQUISITO--EVIDENCIA}

\begin{longtable}{L{1cm}L{4.5cm}L{8cm}}
\caption{Trazabilidad de requisitos críticos de la formulación explícita. Fuente: contrato físico, pruebas y expedientes de ejecución.}\label{tab:explicit-requirements}\\
\toprule ID & Requisito & Evidencia y límite \\
\midrule\endfirsthead
\toprule ID & Requisito & Evidencia y límite \\
\midrule\endhead
R1 & PDE en conductor, cada capa y suelo. & Pruebas parametrizadas de residuos y retropropagación de todas las variantes; campos FEM con etiquetas de cada material. La presencia de un residuo no acredita por sí sola aceptación numérica. \\
R2 & Continuidad de temperatura y flujo. & Pérdidas de transmisión y evaluación independiente de cada interfaz. En la red mixta se comprueban ambos flujos. Los fallos permanecen como rechazos. \\
R3 & Resolución de escalas milimétricas y métricas. & Coordenadas locales, presupuesto por material y pruebas de cobertura; comparación B a 1x, 2x y 4x con FEM. La admisibilidad se separa de la insensibilidad. \\
R4 & Fuente eléctrica trazable y no duplicada. & Contrato DC/AC, pruebas de integral de fuente piel y conservación de corriente DC con temperatura local; comparación FEM de perfiles a igual potencia. La procedencia eléctrica real sigue condicionada a fichas auténticas. \\
R5 & Referencia del mismo problema. & Identidad física, coordenadas comunes, SHA-256 del campo y puertas de refinamiento FEM. Las referencias reducidas no son intercambiables. \\
R6 & Elección sin seleccionar una semilla favorable. & Manifiestos A--C, configuración congelada y D con semillas 71, 83 y 97. Se publican todos los intentos. \\
R7 & Corriente límite con PDE interior. & Raíces FEM de tres mallas, corriente positiva aprendida y nueve intentos de confirmación. El error térmico aislado no sustituye las puertas completas. \\
R8 & Reproducción y documento consistente. & Fuentes y nubes archivadas, repetición A/B del mismo caso y semillas, pruebas automáticas, auditoría de referencias y construcción de tablas y PDF desde los registros. \\
\bottomrule
\end{longtable}

Las pruebas de operadores temporales verifican la forma matemática del término de almacenamiento y la exigencia de capacidad documentada. Una simulación transitoria completa no forma parte de la cobertura ejecutada.

\section{CASOS POR ETAPA Y COBERTURA}
\label{anx:catalogo-casos}

Los casos físicos están en \path{Benchmarks/explicit_study/cases/}. A utiliza coaxial angular y XLPE homogéneo. B y C utilizan XLPE homogéneo y estratos con salto exacto. D conserva esos dos casos para confirmación con semillas nuevas y añade relleno mejorado, zona seca cercana, zona seca lejana, Aras plano y Kim arena.

Los controles analíticos, manufacturados y de piel tienen otra finalidad y no se suman como réplicas de confirmación. La ampacidad DC añade un par heterogéneo--homogéneo con una regla explícita de promedio del suelo. Las geometrías publicadas se adaptan a contornos y fuentes controladas; sus temperaturas no se atribuyen al artefacto. El catálogo más amplio del plan incluye fuentes de propiedades y extensiones no ejecutadas, como humedad acoplada, k(T), datos estacionales y régimen transitorio. La presencia de un caso en archivos no demuestra cobertura experimental.

\section{REPRODUCCIÓN Y CONTRATO DE ARCHIVOS}
\label{anx:protocolo-reproduccion}

Los protocolos registrados son \path{docs/PROTOCOLO_EXPERIMENTAL_EXPLICITO.md} y \path{docs/PROTOCOLO_AMPACIDAD_EXPLICITA.md}, complementados por \path{docs/PROTOCOLO_CALIBRACION_INTERFACES.md}. La enmienda C2 responde a evidencia de B y precede a C y D. La copia congelada de la campaña y los manifiestos por etapa permiten reconstruir su secuencia. Los datos físicos se separan de anchura, profundidad, colocaciones, tasa, presupuesto y semillas. La identidad física incluye PDE, fuentes, interfaces y supuestos eléctricos; un campo reducido no puede sustituir una referencia explícita.

La secuencia comprende: ejecutar pruebas de operadores y muestreo; construir referencias FEM y pasar las puertas de malla; ejecutar A y elegir arquitectura; ejecutar B y analizar colocaciones; ejecutar C y C2 y congelar entrenamiento; ejecutar D con semillas independientes; resolver la extensión de corriente límite; generar tablas y figuras desde los JSON y compilar la tesis. Los comandos principales son \path{Benchmarks/full_study.py}, \path{Benchmarks/full_ampacity_study.py} y \path{Benchmarks/full_study_report.py}. El generador histórico no debe introducir nuevamente cifras reducidas en la tesis activa.

Los NPZ contienen coordenadas, regiones, pesos y temperaturas. Los expedientes PINN incluyen configuración, nubes iniciales y adaptadas, pesos entrenados, estados de optimización, fuentes e historial. Los estados FEM finales incluyen malla curva, etiquetas y grados de libertad. Los ensayos intermedios de la raíz de corriente conservan valores, balances y criterios; no se cuentan como problemas físicos adicionales.

Las semillas 11 y 23 y sus muestreos 101 y 103 pertenecen al desarrollo. Las semillas 71, 83 y 97 y sus muestreos 1001, 1003 y 1007 pertenecen a confirmación. El orquestador reutiliza solo resultados completos de la misma configuración. No se afirma reanudación exacta a mitad de L-BFGS: una ejecución interrumpida debe repetirse desde su configuración o disponer de un protocolo de restauración específico.

La PINN usa Python y PyTorch en Windows, precisión doble y un hilo por proceso; FEniCSx y Gmsh se ejecutan en Linux. Las versiones utilizadas se registran por corrida. Las reinstalaciones y repeticiones del modelo anterior no prueban una reinstalación limpia de este flujo. Las huellas permiten auditar procedencia sin inventar un commit para cambios de trabajo no confirmados en Git.

\section{EXPEDIENTE COMPLETO DE EJECUCIONES}
\label{anx:expediente-evidencia}

La tabla siguiente conserva las ejecuciones completas de A--D, incluidas las rechazadas. La misma información está en \path{Benchmarks/explicit_study/report/runs.csv}, con resúmenes de selección, convergencia y confirmación. El repositorio conserva el respaldo verificado de la tesis anterior; las tablas activas proceden de la formulación explícita.

\begingroup\footnotesize\setlength{\tabcolsep}{3pt}
\begin{longtable}{L{4cm}L{3.5cm}rrr}
\caption{Todas las ejecuciones de la campaña explícita. Fuente: manifiestos por etapa.}\label{tab:explicit-all}\\
\toprule
Etapa / configuración & Caso & Semilla & RMSE, K & Aceptada\\
\midrule\endfirsthead
\toprule
Etapa / configuración & Caso & Semilla & RMSE, K & Aceptada\\
\midrule\endhead
A/fourier\_w16\_d2 & coaxial\_angular & 11 & 0.507 & No\\
A/fourier\_w16\_d2 & coaxial\_angular & 23 & 0.306 & No\\
A/fourier\_w16\_d2 & xlpe\_single & 11 & 5.544 & No\\
A/fourier\_w16\_d2 & xlpe\_single & 23 & 1.672 & No\\
A/local\_w16\_d2 & coaxial\_angular & 11 & 8.475 & No\\
A/local\_w16\_d2 & coaxial\_angular & 23 & 0.244 & No\\
A/local\_w16\_d2 & xlpe\_single & 11 & 0.932 & No\\
A/local\_w16\_d2 & xlpe\_single & 23 & 0.982 & No\\
A/mixed\_w16\_d2 & coaxial\_angular & 11 & 0.146 & No\\
A/mixed\_w16\_d2 & coaxial\_angular & 23 & 0.185 & No\\
A/mixed\_w16\_d2 & xlpe\_single & 11 & 0.911 & No\\
A/mixed\_w16\_d2 & xlpe\_single & 23 & 0.939 & No\\
A/mixed\_w32\_d3 & coaxial\_angular & 11 & 0.302 & No\\
A/mixed\_w32\_d3 & coaxial\_angular & 23 & 0.242 & No\\
A/mixed\_w32\_d3 & xlpe\_single & 11 & 0.894 & No\\
A/mixed\_w32\_d3 & xlpe\_single & 23 & 1.125 & No\\
A/multipole\_w16\_d2 & coaxial\_angular & 11 & 0.069 & No\\
A/multipole\_w16\_d2 & coaxial\_angular & 23 & 0.023 & No\\
A/multipole\_w16\_d2 & xlpe\_single & 11 & 0.381 & No\\
A/multipole\_w16\_d2 & xlpe\_single & 23 & 0.329 & No\\
A/subdomain\_w16\_d2 & coaxial\_angular & 11 & 0.068 & No\\
A/subdomain\_w16\_d2 & coaxial\_angular & 23 & 0.019 & No\\
A/subdomain\_w16\_d2 & xlpe\_single & 11 & 0.357 & Sí\\
A/subdomain\_w16\_d2 & xlpe\_single & 23 & 0.406 & Sí\\
A/subdomain\_w16\_d3 & coaxial\_angular & 11 & 0.032 & No\\
A/subdomain\_w16\_d3 & coaxial\_angular & 23 & 0.059 & No\\
A/subdomain\_w16\_d3 & xlpe\_single & 11 & 0.744 & No\\
A/subdomain\_w16\_d3 & xlpe\_single & 23 & 0.479 & No\\
A/subdomain\_w32\_d3 & coaxial\_angular & 11 & 0.014 & Sí\\
A/subdomain\_w32\_d3 & coaxial\_angular & 23 & 0.017 & No\\
A/subdomain\_w32\_d3 & xlpe\_single & 11 & 0.452 & Sí\\
A/subdomain\_w32\_d3 & xlpe\_single & 23 & 0.372 & Sí\\
A/subdomain\_w8\_d2 & coaxial\_angular & 11 & 0.123 & No\\
A/subdomain\_w8\_d2 & coaxial\_angular & 23 & 0.067 & No\\
A/subdomain\_w8\_d2 & xlpe\_single & 11 & 0.426 & No\\
A/subdomain\_w8\_d2 & xlpe\_single & 23 & 0.820 & No\\
B/interface\_2x & xlpe\_discrete\_layers & 11 & 0.230 & No\\
B/interface\_2x & xlpe\_discrete\_layers & 23 & 0.545 & No\\
B/interface\_2x & xlpe\_single & 11 & 0.294 & Sí\\
B/interface\_2x & xlpe\_single & 23 & 0.618 & Sí\\
B/mixed\_1x & xlpe\_discrete\_layers & 11 & 0.329 & No\\
B/mixed\_1x & xlpe\_discrete\_layers & 23 & 0.347 & No\\
B/mixed\_1x & xlpe\_single & 11 & 0.452 & Sí\\
B/mixed\_1x & xlpe\_single & 23 & 0.372 & Sí\\
B/mixed\_2x & xlpe\_discrete\_layers & 11 & 0.381 & No\\
B/mixed\_2x & xlpe\_discrete\_layers & 23 & 0.235 & No\\
B/mixed\_2x & xlpe\_single & 11 & 0.365 & Sí\\
B/mixed\_2x & xlpe\_single & 23 & 0.329 & Sí\\
B/mixed\_4x & xlpe\_discrete\_layers & 11 & 0.353 & No\\
B/mixed\_4x & xlpe\_discrete\_layers & 23 & 0.218 & No\\
B/mixed\_4x & xlpe\_single & 11 & 0.469 & Sí\\
B/mixed\_4x & xlpe\_single & 23 & 0.591 & Sí\\
B/residual\_2x & xlpe\_discrete\_layers & 11 & 0.315 & No\\
B/residual\_2x & xlpe\_discrete\_layers & 23 & 0.350 & No\\
B/residual\_2x & xlpe\_single & 11 & 0.303 & Sí\\
B/residual\_2x & xlpe\_single & 23 & 0.372 & Sí\\
B/uniform\_2x & xlpe\_discrete\_layers & 11 & 0.490 & No\\
B/uniform\_2x & xlpe\_discrete\_layers & 23 & 0.428 & No\\
B/uniform\_2x & xlpe\_single & 11 & 0.627 & No\\
B/uniform\_2x & xlpe\_single & 23 & 0.249 & No\\
\bottomrule\end{longtable}\endgroup

